\section{Notasi, Prioritas, dan Spesifikasi Token}
\label{sec:ch05-notasi}

\subsection{Maximal Munch}

Ketika beberapa pattern cocok pada posisi yang sama, lexer memilih
lexeme terpanjang. Aturan ini disebut \eng{maximal munch} atau
\eng{longest match}. Contohnya, pada masukan \code{<=}, lexer harus
memilih token \code{LE}, bukan token \code{LT} diikuti
\code{ASSIGN}. Pada masukan \code{==}, lexer memilih \code{EQ},
bukan dua kali \code{ASSIGN}.

\begin{figure}[H]
\centering
\begin{tikzpicture}[node distance=7mm and 14mm]
  \node[phase, text width=3.2cm] (input) {Masukan: \code{a <= b}};
  \node[phase, text width=3.2cm, below left=of input] (salah) {Salah: \code{<} lalu \code{=}};
  \node[phase, text width=3.2cm, below right=of input] (benar) {Benar: \code{<=}};
  \draw[arrow] (input) -- (salah);
  \draw[arrow] (input) -- (benar);
\end{tikzpicture}
\caption{Maximal munch memilih \code{<=} sebagai satu token}
\label{fig:ch05-munch}
\end{figure}

Jika dua pattern menghasilkan panjang yang sama, aturan yang muncul
lebih dahulu dalam spesifikasi menang. Karena itu, urutan aturan
harus dirancang dengan hati-hati. Kata kunci biasanya ditempatkan
sebelum identifier, atau identifier diperiksa terhadap tabel kata
kunci setelah dicocokkan.

\subsection{Notasi di Flex}

Flex memakai notasi yang mirip ekspresi reguler dengan beberapa
tambahan:
\begin{itemize}[leftmargin=*]
  \item \code{[A-Za-z]} menyatakan kelas karakter;
  \item \code{[\^{}0-9]} menyatakan komplemen kelas;
  \item \code{.} mencocokkan karakter apa pun kecuali baris baru;
  \item \code{\textbackslash n}, \code{\textbackslash t} adalah escape;
  \item \code{"..."} mencocokkan string literal;
  \item \code{\{m,n\}} menyatakan pengulangan;
  \item \code{\^{}} dan \code{\$} bergantung pada mode.
\end{itemize}

Flex juga menyediakan \eng{start condition} dan \eng{trailing
context}. Fitur ini memungkinkan lexer menangani mode khusus,
misalnya di dalam komentar atau string. Untuk proyek Subset-C, kita
membatasi diri pada fitur dasar agar lexer tangan dan lexer Flex
dapat dibandingkan secara adil.

\subsection{Aturan Spesifikasi}

Spesifikasi token yang baik memenuhi aturan berikut:
\begin{enumerate}[label=\arabic*.,leftmargin=*]
  \item setiap token memiliki pattern yang jelas;
  \item operator yang lebih panjang ditempatkan sebelum yang lebih
        pendek;
  \item literal float ditempatkan sebelum integer;
  \item komentar dan whitespace ditangani lebih dahulu;
  \item setiap karakter ilegal memiliki perilaku yang terdefinisi;
  \item setiap aturan memiliki contoh dan tes.
\end{enumerate}

\begin{table}[H]
\centering
\caption{Prioritas aturan token Subset-C}
\label{tab:ch05-prioritas}
\begin{tabularx}{\textwidth}{@{}c l X@{}}
\toprule
\textbf{Urutan} & \textbf{Kategori} & \textbf{Contoh} \\
\midrule
1 & whitespace dan komentar & spasi, \code{//}, \code{/* */} \\
2 & operator dua karakter & \code{<=}, \code{>=}, \code{==}, \code{!=}, \code{\&\&}, \code{||} \\
3 & operator satu karakter & \code{+}, \code{-}, \code{*}, \code{/}, \code{<}, \code{=} \\
4 & literal float & \code{3.14}, \code{0.5} \\
5 & literal integer & \code{0}, \code{42} \\
6 & identifier dan kata kunci & \code{hasil}, \code{if} \\
7 & literal char dan string & \code{'a'}, \code{"halo"} \\
8 & karakter ilegal & \code{@}, \code{\#} \\
\bottomrule
\end{tabularx}
\end{table}

\subsection{Kesalahan Leksikal}

Kesalahan leksikal tidak boleh membuat lexer berhenti total. Lexer
harus menghasilkan token kesalahan, mencatat lokasi, lalu melanjutkan
ke karakter berikutnya. Strategi pemulihan sederhana adalah
\eng{skip one character}: buang karakter ilegal dan lanjut.

Untuk string atau komentar yang tidak ditutup, lexer dapat mencari
akhir baris terdekat, menghasilkan token kesalahan, lalu melanjutkan.
Pesan harus menyebut lokasi awal dan jenis masalah, misalnya
\code{baris 7, kolom 3: string tidak ditutup}.
